\scanpage{25}
\mspage{6}
\margin{\uuline{endomorphismes} serait plus exact --- on ne sait pas alors \unsure{exprimer}, en termes de $(p,\beta)$, que $u$ est un \uuline{automorphisme}. \unsure{Ça revient à un|Ça n'est pas un} endomorphisme, \unsure{imposant} seulement $p\in\widehat{\mathbb Z}$, non $p\in\widehat{\mathbb Z}^{*}$, en tous cas $p$ de la forme $2\gamma+1$, avec $\gamma\in\widehat{\mathbb Z}$.}
$u$ autom \inserted{du groupoïde $\widehat{\mathrm{II}}_J$} qui commute à $\unsure{\mathcal G_J|G_J}$, et fixe les pts,
\deleted{\unsure{fini et déterminé}}
et qui \unsure{invarie} \deleted{\unsure{tous}} les sous-groupoïdes $\widehat{\mathrm{II}}_{J,i}$, où $\widehat{\mathrm{II}}_{J,i}$ est engendré par $R_i^\omega$, $a_i^\omega$ \unsure{$\in$}, $\omega\in\uline{\omega}$. La condition : $\mathrm{II}_{J,i}$ \deleted{défini} donné par
\[
u(a_i^\omega)=a_i^\omega\underbrace{\varphi_{i\infty}^\omega(l_0^\gamma)}_{(\lambda_i^\omega)^\gamma}
=\varphi_i^{\omega'}(l_\infty^\gamma)a_i^\omega
\qquad\left|\begin{array}{l}l_0^\gamma\text{ indépendant de }i,\omega\\\gamma\in\widehat{\mathbb Z}\end{array}\right.
\]
on aura
\[
u(\deleted{\unsure{a}}a_i^{\omega'}a_i^\omega)=\varphi_{i\infty}(l_0^{2\gamma}\deleted{\unsure{1}}l_0)
\]
\[
u\bigl(\underbrace{\unsure{\varphi_{i\infty}^{\omega}|\varphi_{i\infty}^{\omega'}}(l_0)}_{\lambda_i^\omega}\bigr)
=\underbrace{\varphi_i^\omega(l_0^p)}_{(\lambda_i^\omega)^p}
\qquad\gamma=\frac{p-1}{2}
\]
\margin{faire les calculs en partant d'abord de $u(a_i^\omega)=a_i^\omega\varphi_i^\omega(\lambda)$ avec $\lambda\in\widehat\Pi_{0,3}$ qlc., (et non $\lambda$ de la forme $l_0^\gamma$, $\gamma\in\widehat{\mathbb Z}\ldots$)}
\[
\left\{\begin{aligned}
u(a_i^\omega)&=a_i^\omega\varphi_{i\infty}(l_0^\gamma)
&\gamma&=\frac{p-1}{2}\quad p\in\widehat{\mathbb Z}^{*}\text{ multiplicateur}\\
u(b_i^\omega)&=b_i^\omega\varphi_{\omega i}^\omega(\deleted{\unsure{\xi}}\beta)
&\beta&\in\widehat\Pi_{0,3}\\
&\deleted{\unsure{\varphi_{\omega'i}^{\omega'}(\beta)b_i^\omega}}\quad
\inserted{\bigl(=\varphi_{\omega'i}^{\omega'}(\tau_\infty\beta)b_i^\omega\bigr)}
\end{aligned}\right.
\]
Ici \deleted{\unsure{le}} $(p,\beta)$ détermine $u$ (en fait, $\beta$ détermine $p\ldots$)
\[
u(\widetilde b_{\unsure{t}i}^\omega)=\widetilde b_i^\omega\varphi_{\omega i}^\omega
\bigl(\underbrace{l_1^{-\gamma}\beta}_{\unsure{\sigma_\infty(\lambda)}\beta}\bigr)
=\varphi_{\omega'i}^{\omega'}\bigl(l_0^{-\gamma}\rho^{-1}(\beta)\bigr)\widetilde b_i^\omega
\]
Les équations sur $(p,\beta)$ équivalent à \unsure{sont|$\longrightarrow$}
\[
(12)\quad\left\{\begin{aligned}
\beta\sigma_\infty(\beta)&=1\\
\alpha\rho^{-1}(\alpha)\rho^{-2}(\alpha)&=1\\
\alpha&=\underbrace{l_1^{-\gamma}}_{\sigma_\infty(\lambda)}\beta
\end{aligned}\right.
\qquad\deleted{\text{où }\alpha=l_1^{-\gamma}\beta}\quad\text{ou}\quad
\bigl(\text{ou }\alpha^{-1}\rho(\alpha^{-1})\rho^2(\alpha^{-1})=1\bigr)
\]
\margin{\uline{NB} Si on remplace $\beta$ par $\beta^{-1}$, $\alpha$ par $\alpha^{-1}$ on retrouve les équations habituelles. cf plus bas}
\margin{Si on change la définition de $\beta$, en faisant que $\beta$ devienne $\sigma_\infty(\beta)$, les deux premières équations (12) restent pareilles mais $\alpha=l_1^{\unsure{\gamma}}\sigma_\infty(\beta)$, ce qui est un peu moins joli\ldots\ et est \unsure{évident} \unsure{que} $\beta\mapsto\beta'$.}
Notons que si $u$ est défini par \inserted{$(p,$}\deleted{$\beta$}\inserted{$\beta)$}, alors
$\tau(u)$ \inserted{$=\tau u\tau^{-1}$} est défini par $(p,\beta'=\tau_\infty(\beta))$. Il faudrait pouvoir vérifier directement, par un calcul dans $\Pi_{0,3}$, que les équations (12) en $\beta$, impliquent les $\Rightarrow$ relations pour $\beta'=\tau_\infty(\beta)$ (avec $\unsure{\gamma\prime|\gamma}=\gamma$).
\margin{24}
